\abstrakindonesia{Rumah Makan Nasi Gerilya merupakan usaha nasi Padang di Kota Medan yang menggunakan GrabFood sebagai salah satu kanal penjualan utama, tetapi kondisi di lokasi penelitian menunjukkan penjualan yang belum stabil, pembelian ulang yang menurun, adanya keluhan pelanggan, serta tekanan persaingan dari merchant sejenis pada platform yang sama. Tujuan penelitian ini adalah mengidentifikasi faktor internal dan eksternal Rumah Makan Nasi Gerilya pada platform GrabFood serta merumuskan strategi pemasaran berdasarkan analisis SWOT untuk mendukung peningkatan penjualan. Penelitian dilakukan pada Rumah Makan Nasi Gerilya di Kota Medan dengan fokus pada aktivitas pemasaran melalui GrabFood. Pengumpulan data dilakukan melalui wawancara dengan pemilik, pegawai, pelanggan, dan informan pembanding, observasi platform GrabFood, wawancara pelanggan, dokumentasi usaha, ulasan pelanggan, serta studi pustaka. Hasil penelitian menunjukkan bahwa skor IFAS sebesar 2,50 dan skor EFAS sebesar 2,51 menempatkan Nasi Gerilya pada kondisi sedang dengan posisi Kuadran I, yaitu strategi agresif atau SO karena kekuatan internal dan peluang eksternal lebih dominan daripada kelemahan dan ancaman. Kekuatan utama usaha terletak pada cita rasa lauk, porsi besar, reputasi digital, alur pengecekan pesanan, dan respons keluhan, sedangkan kelemahan utama terdapat pada akurasi packing, antrean lintas kanal, update stok, persepsi harga, dan kejelasan deskripsi paket. Strategi alternatif yang dirumuskan meliputi strategi SO, WO, ST, dan WT, dengan prioritas pada perbaikan sistem antrean dan checklist packing, penguatan paket menu unggulan berbasis segmen, perbaikan tampilan digital dan menu digital, sinkronisasi promo dengan stok, kapasitas, HPP, dan margin, serta pemanfaatan stamp digital, database pelanggan, dan Flowie untuk mendorong pembelian ulang.}

\katakunciindonesia{strategi pemasaran, GrabFood, SWOT, UMKM kuliner, penjualan}
